Gabe's paper (aka math formalism) --- DONE.

Paper B --- 2D Li6 segmented detectors (applying the formalism and comparing for "conventional" paper A / Chooz formalism --- In Review.

Geo-neutrino paper (might be part of Whitepaper) --- finishing

directionality for continuous sources (unlike reactors) 


Paper C --- AI/ML paper (Max)
Other formalism (KS, Kuiper / circular stats, maybe) .....
Adding cut on capture timing --- potential to improve directionality (can work with existing simulations, potentially low-hanging fruit).
Adding cuts in general.

SMR paper --- how it all works for multi-core (recall the problem in 2007 paper about Chooz)

Future fusion-fission hybrid reactors (antineutrinos will be from the fission blanket but the geometry might play a role for a particular design). Discuss how antineutrino detection might play a role in future safety and safeguards of fusion-fission reactors.

Other use cases --- spent fuel? (DoubleChooz PRL first measurement)

Frost paper (hardware test with collimated gamma-ray source, software: tuo or mtc-b RATPAC2 simulations)

Non-square segments: mathematical formalism should be the same? Hexagonal, triangle.

Adding different weights, e.g. central segment is less important than adjacent segments, currently the CFND formalism does not have weights.

Paper to deal with realistic backgrounds --- does directionality help with rejecting backgrounds? How backgrounds play a role? Also include optimal bin size for best directionality.


SiPM and scintillator under high-pressure (UH high-pressure test stand) --- on hold. need funding

Explore SNO+ data (table 1, arXiv:2505.04469) --- captures on hydrogen, do RAT-PAC simulations. 
Geo: continuous. Reactor: distributed (reactors with known positions).
Can we get the same raw-reconstructed data for KamLAND and Borexino?

glass / plastic checker-board --- Li-6 doped glass idea

IMB old data recovery  (from Bob)

Bugey old data recovery (from original Bugey team)
